% GENERATED FILE. Edit canonical lessons, then run npm run generate:books.
% canonical-source-hash: fnv1a64:a9d5ef6c3cc34c3d
% canonical-lessons: PT-C406-avaria, PT-C406-avariar, PT-C406-rebocar, PT-C406-reboque, PT-C406-oficina

\chapter{A avaria}
\label{ch:pt-a-avaria}
\hlchaptermodality{406}

\begin{chapteropening}
\textbf{By the end of this chapter:} \emph{I can deal with a breakdown: the breakdown, breaking down, towing, the tow truck and the repair garage.}

Complete the last lesson of chapter 406: I can deal with a breakdown: the breakdown, breaking down, towing, the tow truck and the repair garage.
\end{chapteropening}

\section[a avaria]{a avaria — the breakdown, the fault}

\label{lesson:PT-C406-avaria}



\begin{quote}
Before the new one: say the Portuguese for the wheel rim, then the Portuguese for the warning triangle.
\end{quote}

\subsection*{You'll want to know: a avaria}
\textbf{a avaria} — \textquotedblleft{}the breakdown, the fault\textquotedblright{}.

\begin{quote}
O carro teve uma avaria a meio da viagem.
\end{quote}

\begin{quote}
\emph{The car broke down halfway through the trip.}
\end{quote}

\emph{Estar avariado} — to be out of order. In Brazil: \emph{o defeito}, \emph{a pane}.

\begin{cousinweb}
\emph{avaria} came through Italian \emph{avaria} from Arabic \emph{ʿawāriyya} 'damaged goods', built on \emph{ʿawār} 'fault, flaw' (EtymDB). It began as a word for damage to a ship's cargo.
\end{cousinweb}

\begin{grammarlens}[title={What you've built}]
One more word for telling, explaining and describing at B1.
\end{grammarlens}

\subsection*{Your turn}
\begin{itemize}
\raggedright
  \item \emph{Say it:} \emph{a avaria}
  \item \emph{Say it:} \emph{a avaria}, once more
  \item \emph{Say it:} a sentence of your own with \emph{a avaria}
  \item \emph{Recall:} say the Portuguese for the wheel rim, then the Portuguese for the warning triangle, then say \emph{a avaria} again
\end{itemize}

\begin{tcolorbox}[breakable,colback=teal!4,colframe=teal!35!black,title={Before you move on}]
What does a avaria mean? (\textquotedblleft{}The breakdown, the fault\textquotedblright{}.) Say it once more.
\end{tcolorbox}
\section[avariar]{avariar — to break down}

\label{lesson:PT-C406-avariar}



\begin{quote}
Before the new one: say the Portuguese for the warning triangle, then the Portuguese for the breakdown.
\end{quote}

\subsection*{You'll want to know: avariar}
\textbf{avariar} — \textquotedblleft{}to break down\textquotedblright{}.

\begin{quote}
A máquina de café avariou outra vez.
\end{quote}

\begin{quote}
\emph{The coffee machine broke down again.}
\end{quote}

Mostly for machines: \emph{o elevador avariou}.

\begin{cousinweb}
\emph{avariar} is built on \emph{avaria} 'fault', which came through Italian from Arabic \emph{ʿawāriyya} 'damaged goods' (EtymDB).
\end{cousinweb}

\begin{grammarlens}[title={What you've built}]
One more word for telling, explaining and describing at B1.
\end{grammarlens}

\subsection*{Your turn}
\begin{itemize}
\raggedright
  \item \emph{Say it:} \emph{avariar}
  \item \emph{Say it:} \emph{avariar}, once more
  \item \emph{Say it:} a sentence of your own with \emph{avariar}
  \item \emph{Recall:} say the Portuguese for the warning triangle, then the Portuguese for the breakdown, then say \emph{avariar} again
\end{itemize}

\begin{tcolorbox}[breakable,colback=teal!4,colframe=teal!35!black,title={Before you move on}]
What does avariar mean? (\textquotedblleft{}To break down\textquotedblright{}.) Say it once more.
\end{tcolorbox}
\section[rebocar]{rebocar — to tow}

\label{lesson:PT-C406-rebocar}



\begin{quote}
Before the new one: say the Portuguese for the breakdown, then the Portuguese for to break down.
\end{quote}

\subsection*{You'll want to know: rebocar}
\textbf{rebocar} — \textquotedblleft{}to tow\textquotedblright{}.

\begin{quote}
Um camião rebocou o carro até à cidade.
\end{quote}

\begin{quote}
\emph{A truck towed the car to the town.}
\end{quote}

\emph{Rebocar} also means to plaster a wall, a different word.

\begin{cousinweb}
\emph{rebocar} 'to tow' is inherited from Vulgar Latin \emph{remulcāre}, built on Latin \emph{remulcum} 'towrope' (EtymDB).
\end{cousinweb}

\begin{grammarlens}[title={What you've built}]
One more word for telling, explaining and describing at B1.
\end{grammarlens}

\subsection*{Your turn}
\begin{itemize}
\raggedright
  \item \emph{Say it:} \emph{rebocar}
  \item \emph{Say it:} \emph{rebocar}, once more
  \item \emph{Say it:} a sentence of your own with \emph{rebocar}
  \item \emph{Recall:} say the Portuguese for the breakdown, then the Portuguese for to break down, then say \emph{rebocar} again
\end{itemize}

\begin{tcolorbox}[breakable,colback=teal!4,colframe=teal!35!black,title={Before you move on}]
What does rebocar mean? (\textquotedblleft{}To tow\textquotedblright{}.) Say it once more.
\end{tcolorbox}
\section[o reboque]{o reboque — the tow truck; the tow, the trailer}

\label{lesson:PT-C406-reboque}



\begin{quote}
Before the new one: say the Portuguese for to break down, then the Portuguese for to tow.
\end{quote}

\subsection*{You'll want to know: o reboque}
\textbf{o reboque} — \textquotedblleft{}the tow truck; the tow, the trailer\textquotedblright{}.

\begin{quote}
Chamámos o reboque e esperámos uma hora.
\end{quote}

\begin{quote}
\emph{We called the tow truck and waited an hour.}
\end{quote}

In Brazil: \emph{o guincho}.

\begin{cousinweb}
\emph{reboque} is built on \emph{rebocar} 'to tow', which goes back to Latin \emph{remulcum} 'towrope' (EtymDB).
\end{cousinweb}

\begin{grammarlens}[title={What you've built}]
One more word for telling, explaining and describing at B1.
\end{grammarlens}

\subsection*{Your turn}
\begin{itemize}
\raggedright
  \item \emph{Say it:} \emph{o reboque}
  \item \emph{Say it:} \emph{o reboque}, once more
  \item \emph{Say it:} a sentence of your own with \emph{o reboque}
  \item \emph{Recall:} say the Portuguese for to break down, then the Portuguese for to tow, then say \emph{o reboque} again
\end{itemize}

\begin{tcolorbox}[breakable,colback=teal!4,colframe=teal!35!black,title={Before you move on}]
What does o reboque mean? (\textquotedblleft{}The tow truck; the tow, the trailer\textquotedblright{}.) Say it once more.
\end{tcolorbox}
\section[a oficina]{a oficina — the repair garage; the workshop}

\label{lesson:PT-C406-oficina}



\begin{quote}
Before the new one: say the Portuguese for to tow, then the Portuguese for the tow truck.
\end{quote}

\subsection*{You'll want to know: a oficina}
\textbf{a oficina} — \textquotedblleft{}the repair garage; the workshop\textquotedblright{}.

\begin{quote}
O mecânico da oficina disse que demora dois dias.
\end{quote}

\begin{quote}
\emph{The mechanic at the garage said it takes two days.}
\end{quote}

A garage to park in is \emph{a garagem}; repairs happen at \emph{a oficina}.

\begin{cousinweb}
\emph{oficina} is a learned borrowing from Latin \emph{officīna} 'workshop' (EtymDB), related to \emph{opifex} 'worker', built on \emph{opus} 'work'.
\end{cousinweb}

\begin{grammarlens}[title={What you've built}]
One more word for telling, explaining and describing at B1.
\end{grammarlens}

\subsection*{Your turn}
\begin{itemize}
\raggedright
  \item \emph{Say it:} \emph{a oficina}
  \item \emph{Say it:} \emph{a oficina}, once more
  \item \emph{Say it:} a sentence of your own with \emph{a oficina}
  \item \emph{Recall:} say the Portuguese for to tow, then the Portuguese for the tow truck, then say \emph{a oficina} again
\end{itemize}

\begin{tcolorbox}[breakable,colback=teal!4,colframe=teal!35!black,title={Before you move on}]
What does a oficina mean? (\textquotedblleft{}The repair garage; the workshop\textquotedblright{}.) Say it once more.
\end{tcolorbox}
